% Chapter 13 — Frontend Interaction (Data Perspective)

\chapter{Frontend Interaction (Data Perspective)}
\label{ch:frontend}

\section{Browsers never run SQL}
The React application only speaks HTTP JSON to Express. Students never receive raw database passwords or service keys. That separation is fundamental: \textbf{the API is the gatekeeper}.

\section{Typical data round-trip}
\begin{enumerate}[leftmargin=*]
  \item User logs in $\rightarrow$ receives JWT stored locally (e.g.\ \texttt{localStorage}).

  \item Axios attaches \texttt{Authorization: Bearer <token>} on later requests.

  \item A page like course detail loads course JSON from \texttt{GET /courses/:id}. That response already joins instructor display names via nested selects on the server side.

  \item Video tracking sends periodic \texttt{POST /activity} and \texttt{PUT /progress/...} calls---these become inserts/upserts in PostgreSQL.

\end{enumerate}

\section{Quiz authoring note}
Instructors may import question banks from JSON in the UI; each question becomes repeated API posts (\texttt{POST /quizzes/:id/questions}). That is many small transactions rather than one giant bulk insert---easier for coursework debugging.

\section{Contract guarantees (what we promise)}
\begin{itemize}[leftmargin=*]
  \item Final quiz scores are computed on the server, not trusted from the browser.

  \item Dashboard numbers come from API JSON responses (often fed by materialized views), not from guesswork in React state.

  \item On HTTP 401, the client clears tokens and redirects to login so half-expired sessions do not silently break writes.

\end{itemize}

\section{Failure handling}
Network timeouts and validation errors surface as toast messages or inline errors. From a DBMS angle, the important lesson is: \textbf{failed requests mean no row was committed}---the database stays consistent even when Wi-Fi drops mid-form.
